% ============================
% ME 793 – Computational Fluid Dynamics
% Project 2 Report Template
% REVTeX 4.2e (revtex4-2) + AAPM class
% ============================
\documentclass[
  aapm,
  amsmath,amssymb,
  reprint,          % two-column journal-like layout
  %preprint,        % uncomment for single-column draft
  floatfix
]{revtex4-2}

% ----------------------------
% Packages
% ----------------------------
\usepackage{graphicx}       % figures
\usepackage{dcolumn}        % decimal-aligned table columns
\usepackage{bm}             % bold math symbols
\usepackage{siunitx}        % SI units (optional, but useful)
\usepackage{hyperref}       % clickable references/links
\usepackage{algorithm2e}    % pseudocode (optional)
\usepackage{tabularx}

% ----------------------------
% Title block
% ----------------------------
\begin{document}

\title{Project 3 Proposal: Flow Around an Actuator Disk Model for UAV Wind Sensor Placement}

\author{Jaden Mecham}
\email{jadenmecham@unr.edu}
\affiliation{Department of Mechanical Engineering, University of Nevada, Reno}

\date{\today}

\maketitle

\section{Problem Description}

In multirotor drone design, accurate wind sensing is essential for flight stabilization
and environmental monitoring, but propeller-induced flow can corrupt measurements
from onboard wind sensors.  The present study simulates the two-dimensional incompressible
flow generated by a hovering rotor using an actuator-disk model. The propeller is modeled
by a thin band of body force that injects momentum into the fluid.
The flow is governed by the incompressible Navier--Stokes equations,
\begin{equation}
  \frac{\partial \mathbf{u}}{\partial t}
  + (\mathbf{u}\cdot\nabla)\mathbf{u}
  = -\nabla p + \frac{1}{Re}\,\nabla^{2}\mathbf{u} + \mathbf{f},
  \qquad \nabla\cdot\mathbf{u} = 0,
\end{equation}
where $\mathbf{f}$ is the actuator-disk body force and
$Re = v_{h} D / \nu$ is the Reynolds number formed from the induced velocity
$v_{h} = \sqrt{T/(2\rho A)}$, the disk diameter $D$, and the kinematic viscosity $\nu$.
Here $T$ is the rotor thrust and $A = \pi D^{2}/4$ is the disk area.
Target Reynolds numbers of $Re \in [100,\,1000]$ will be explored.

\section{Numerical Approach}

The solver extends the fractional-step (projection) method developed in Project~2
on a staggered grid.  The only structural modification is the addition of the
actuator-disk body force $\mathbf{f}$ to the momentum predictor step:
\begin{equation}
  \mathbf{u}^{*} = \mathbf{u}^{n}
    + \Delta t\!\left[-A(\mathbf{u}^{n})
      + \frac{1}{Re}\,L\,\mathbf{u}^{n} + \mathbf{f}\right],
\end{equation}
where $A$ is the discrete advection operator and $L$ is the vector Laplacian.
The disk force is applied uniformly over a single column of $u$-face cells spanning
the rotor diameter, yielding a spatially compact forcing that is straightforward to
implement within the existing sparse-matrix framework.  The pressure correction step
and divergence-free projection remain identical to Project~2.
Second-order central differences are used on a uniform Cartesian grid with explicit
Euler time integration ($CFL \leq 0.5$), and a scipy CG solver handles the pressure
Poisson system at each step.

\section{Planned Features and AI Integration}

\textbf{Actuator-disk body force module.}
    Implement a sparse forcing vector that injects a prescribed thrust coefficient
    $C_T = T/(\tfrac{1}{2}\rho v_{h}^{2} A)$ into the $u$-momentum equation across
    the disk region, with AI assistance to verify the index mapping
    from the 2D staggered grid to the flattened 1D system.

\textbf{Wake velocity post-processing pipeline.}
    Build routines to extract centerline velocity profiles, compute thrust from
    integrated momentum flux, and compare against theory, with AI assistance for
    the quadrature and plotting.

\textbf{Grid-convergence and time-step study.}
    Automate a sweep over grid resolutions ($32^{2}$ to $256^{2}$), log the
    $L^{2}$ velocity error relative to the analytical hover solution, and produce
    convergence-rate plots, using AI to scaffold the sweep and curve-fitting.

\textbf{Sensor-placement analysis.}
    Map $|\mathbf{u} - \mathbf{u}_{\infty}|$ across the domain to identify
    regions of low flow disturbance suitable for anemometer placement, with
    AI assistance for contouring and thresholding.


\section{Validation Plan}

The primary benchmark is Froude--Rankine actuator-disk momentum theory, which predicts
that the streamwise velocity at the disk plane equals the hover-induced velocity
$v_{h}$, and that the far-wake velocity approaches $2v_{h}$.
The simulation will be validated by comparing (i) the centerline
velocity one diameter downstream against the theoretical value $2v_{h}$,
(ii) the integrated thrust extracted from the momentum flux against the prescribed
forcing, and (iii) the radial velocity profile in the far wake against the
top-hat actuator-disk prediction.

\section{Expected Outcomes and Timeline}

\begin{itemize}
  \item \textbf{Week 1 (Apr.\ 24 -- May 1):}
    Integrate body force term into Project~2 solver. verify steady hover solution
    for a single test case at $Re = 200$.Quantitative validation against momentum theory. grid- and time-step convergence
    study. confirm second-order spatial accuracy.
  \item \textbf{Week 2 (May 1 -- May 6):}
    Parameter study over $Re \in [100,\,1000]$; sensor-placement post-processing.
    generate all figures. Draft final report. compile AI Usage Log. Submit.
\end{itemize}


% ============================
% References
% ============================
\bibliographystyle{aapmrev4-2}
\bibliography{references}


\end{document}
